\chapter{The proof system, as deployed and as planned}\label{ch:map}

This chapter is the map. It says what leanVM's proof system does at the pin, interaction by
interaction, where the blueprint puts each part, and where the code that exists sits. The
detailed tables of the transcript, of the verifier's checks and of the disagreements between
the sources, one set per phase, are in \cref{ch:faithfulness}; this chapter gives the shape.

\section{Two fields, one committed column}\label{sec:map-fields}

Everything the prover commits to is a table of values in the field
\begin{define}{The fields $\K$ and $\E$}
$\K = \mathrm{GF}(2^{64})$ is the field of the columns: every cell of every table of the
arithmetization is an element of $\K$, represented as 64 bits. $\E = \mathrm{GF}(2^{192})$ is
the field of the challenges: $\E = \K[y]/(y^3 + y + 1)$, so an element of $\E$ is three
\emph{limbs} in $\K$, and $\K$ sits inside $\E$ as the elements whose upper two limbs are zero.
Both have characteristic 2: $x + x = 0$, so subtraction is addition and $-1 = 1$. In Lean,
CompPoly's \lean{BF64} and \lean{BF64.Ext3}; the review confirmed that their defining
polynomials, bit encodings and limb order are those of the Rust and of the specification
(\cref{ch:libraries}).
\end{define}
The prover's whole witness is one column: every table of the arithmetization (the six opcode
tables, the three memory limbs, the two finalize counts, and the packed bit witness of the
BLAKE2s circuit) is stacked into a single table $q$ of $2^{\mu}$ cells of $\K$, with
$15 \le \mu \le 28$. The verifier never sees $q$; it holds a commitment to it and, at the very
end, asks the commitment scheme to certify weighted sums $\sum_w W(w)\, \widetilde{q}(w)$ over
the cube, where $\widetilde{q}$ is the multilinear extension of $q$ and $W$ is a weight the
verifier can evaluate itself. An evaluation claim $\widetilde{q}(r) = v$ at a point
$r \in \E^{\mu}$ is the case $W = \eq(r, \cdot)$.

\section{The protocol at the pin, phase by phase}\label{sec:map-protocol}

\Cref{fig:phases} shows the six phases the blueprint distinguishes, with what each phase adds
to the \emph{claim pool}: the list of evaluation claims on columns of $q$ that the opening
discharges at the end. The counts below are those of the deployed verifier at the pin
(\file{crates/lean_vm/src/cpu/mod.rs:711-779} is the top-level flow); each is cited in
\cref{ch:faithfulness}.

\paragraph{Setup (before any phase).} The Fiat--Shamir chain is seeded with a digest of the
program and of the Flock circuit's identifier, then with the two public words. The prover
announces eight sizes: the memory log-size $\kappa_{\mathrm{mem}}$, the six tables' log-heights
$\tau_0, \dots, \tau_5$, and the rate exponent $\rho$ of the commitment scheme. The verifier
checks that each is canonically encoded (upper limbs zero), that each public word has a zero
top limb, the caps ($2^{\kappa_{\mathrm{bc}}}$ instructions with $\kappa_{\mathrm{bc}} \le 32$;
$16 \le \kappa_{\mathrm{mem}} \le 32$; $\tau_j \le 32$; $\tau_{\mathrm{BLAKE2S}} \ge 3$;
$1 \le \rho \le 4$), derives every layout, and checks the stacked size
$15 \le \mu \le 28$. None of this is a phase of the oracle protocol; all of it is a check of
the deployed verifier, and \cref{ch:drift} counts it.

\paragraph{Commit.} The prover sends the commitment to $q$: a Merkle root, two field elements
with zero top limbs. In the oracle protocol this is the one oracle message: from here on the
verifier may query $\widetilde{q}$ at points of $\E^{\mu}$ and nothing else.

\paragraph{Bus.} The bus argument shows that the tuples every table \emph{pushes} onto the
three buses (state, memory, bytecode) are, as a multiset, the tuples that are \emph{pulled},
and that no read count is zero. The verifier draws the fingerprint challenge
$\alpha \in \E^4$ and the multiset challenge $\beta \in \E$; each 16-coordinate tuple $t$
becomes the scalar $\beta + \sum_i \eq(\alpha, i)\, t_i$; the pushed scalars, the pulled scalars
and the count cells are each stacked into a tree of $2^{\mu_{\mathrm{bus}}}$ leaves (padded
with $1$), and the prover sends two roots: $R$, claimed to be the product of the pushed leaves
\emph{and} of the pulled leaves, one scalar for both sides, and $R_c$, the product of the
counts. A batched GKR grand-product argument then walks the three trees from the roots to the
leaves, radix 4 (a binary first layer when $\mu_{\mathrm{bus}}$ is odd): per layer, a
combiner $\lambda$ batches the three trees; a sumcheck of $\mu_{\mathrm{bus}} - \ell$ rounds,
each round a cofactor of degree 4 sent as four coefficients; the twelve children; one layer
check; two combination challenges. A fresh combiner is drawn after every layer, the last one
included, where it is used by nothing. The GKR ends with three claims on the leaf tables at one
point $\zeta$ assembled from the last layer's challenges. The verifier checks $R_c \neq 0$,
computes itself the public part of each leaf claim (the selectors, the constants, the index
column, the program's bytecode columns), reads five boundary evaluations the prover sends (the
three memory limbs and the memory finalize count at $\zeta_{< \kappa_{\mathrm{mem}}}$, the
bytecode finalize count at $\zeta_{< \kappa_{\mathrm{bc}}}$), and derives what the tables owe
each side: three values $\mathit{rem}_s$. Into the pool: five column claims.

\paragraph{Table sumcheck.} One sumcheck proves the two constraints of the JUMP table and the
tables' share of the three bus sides together. The verifier draws $\xi \in \E$; the constraints
take $\xi^0, \xi^1$ and the three sides $\xi^2, \xi^3, \xi^4$, shared by every table; the
target is $\sum_s \xi^{2+s}\, \mathit{rem}_s$, derived, never sent. The sumcheck runs
$\tau_{\max}$ rounds, the maximum log-height over the six opcode tables, binding the highest
variable first, each table joining when its height is reached; a round message is three
coefficients of a cubic, the fourth derived from the running claim, so that no round check can
fail. The prover then sends one value per column of every opcode table, 104 in all (the
eighteen value limbs of the BLAKE2S table included), the verifier recomputes the summand from
them and checks it against the running claim. Into the pool: 104 column claims, each at the
prefix $\rho_{< \tau_j}$ of the sumcheck's point; the eighteen limb claims are claims on
strided slots of the packed bit witness.

\paragraph{Public input.} The verifier draws $r \in \E$; the prover sends $c_0, c_1$, claimed
to be the values of the two low memory limbs at $(r, 0, \dots, 0)$, the line through cells
0 and 1. The specification checks each against the public words' limb,
$c_\ell = (1 + r)\, w_{0,\ell} + r\, w_{1,\ell}$; the three executable verifiers (Rust, Python,
the recursion guest) check one combined equation $c_0 + y\, c_1 = (1 + r)\, w_0 + r\, w_1$ in
$\E$. Into the pool: three claims, on the three memory limbs at $(r, 0, \dots, 0)$, with the
values sent and $0$ for the top limb, for which nothing is sent.

\paragraph{BLAKE2s validity (Flock) and ring switching.} Flock proves that the bits packed into
the region $\qflock$ of $q$ satisfy the R1CS of the BLAKE2s compression, block by block, with
position 512 of every block equal to $1$. It commits nothing of its own and reads no claim of
the pool: it is a zerocheck with a univariate skip (64 values, one challenge, $8 + k$ rounds of
two coefficients, two final values) followed by a lincheck ($8$ rounds, then 64 values
$s_0, \dots, s_{63}$), with a single check that can reject, the lincheck's terminal identity.
Ring switching then turns the 64 claims on the bit polynomials into one weighted claim on the
$\qflock$ region of $q$, with six challenges and no message. Into the pool: one weighted claim.

\paragraph{Opening.} The verifier draws $\lambda \in \E$; the ring-switched claim takes
$\lambda^0$ and the 112 pooled point claims the next powers, each point claim first turned
into a weighted claim through the stack's selectors; the batch is one weighted claim
$\sum_w W_\lambda(w)\, \widetilde{q}(w) = C_\lambda$, and the commitment scheme (WHIR over
Reed--Solomon codes of $\K$ in the novel polynomial basis, with Merkle trees over BLAKE2s)
opens it. The specification is explicit: \enquote{This one opening discharges every pooled
claim; there is no separate reduction sumcheck}
(\file{doc/leanvm/body/08-end-to-end-protocol.tex:99}). The verifier finally checks that the
proof stream has been consumed entirely.

\begin{figure}[p]
\centering
\begin{tikzpicture}[
  font=\small\sffamily, node distance=6mm,
  ph/.style={draw, rounded corners=2pt, fill=ok!5, draw=ok!70!black, text width=0.30\linewidth,
    align=left, inner sep=4pt},
  pool/.style={draw, rounded corners=2pt, fill=minor!5, draw=minor!70!black,
    text width=0.24\linewidth, align=left, inner sep=4pt, font=\footnotesize\sffamily},
  chal/.style={text width=0.30\linewidth, align=left, font=\footnotesize\sffamily},
  arr/.style={-{Stealth[length=2mm]}, thick}]
\node[ph] (setup) {\textbf{Setup} \\ seed; 8 announced sizes; the caps; the stacking window};
\node[ph, below=of setup] (commit) {\textbf{Commit} \\ the root of $q$ (2 elements)};
\node[ph, below=of commit] (bus) {\textbf{Bus} \\ $\alpha, \beta$; roots $R, R_c$; batched
  GKR over three trees; $R_c \neq 0$; five boundary evaluations; $\mathit{rem}_s$ derived};
\node[ph, below=of bus] (table) {\textbf{Table sumcheck} \\ $\xi$; $\tau_{\max}$ rounds of 3
  coefficients; 104 column values; the final check};
\node[ph, below=of table] (pub) {\textbf{Public input} \\ $r$; $c_0, c_1$; the line check
  (per limb, or combined)};
\node[ph, below=of pub] (flock) {\textbf{BLAKE2s validity} \\ zerocheck with univariate skip;
  lincheck; the terminal identity; ring switching (6 challenges)};
\node[ph, below=of flock] (open) {\textbf{Opening} \\ $\lambda$; 113 claims batched; WHIR
  over Merkle trees; the stream fully consumed};
\draw[arr] (setup) -- (commit); \draw[arr] (commit) -- (bus); \draw[arr] (bus) -- (table);
\draw[arr] (table) -- (pub); \draw[arr] (pub) -- (flock); \draw[arr] (flock) -- (open);
\node[pool, right=12mm of bus] (p1) {into the pool: 5 column claims at $\zeta$ (the memory
  limbs, the two finalize counts)};
\node[pool, right=12mm of table] (p2) {into the pool: 104 column claims at $\rho_{<\tau_j}$,
  18 of them strided slots of $\qflock$};
\node[pool, right=12mm of pub] (p3) {into the pool: 3 claims on $\mathit{mem}_0, \mathit{mem}_1,
  \mathit{mem}_2$ at $(r, 0, \dots, 0)$};
\node[pool, right=12mm of flock] (p4) {into the pool: 1 weighted claim on the $\qflock$
  region};
\node[pool, right=12mm of open] (p5) {out of the pool: all 113, in one opening};
\draw[arr, minor!70!black] (bus.east) -- (p1.west); \draw[arr, minor!70!black] (table.east) -- (p2.west);
\draw[arr, minor!70!black] (pub.east) -- (p3.west); \draw[arr, minor!70!black] (flock.east) -- (p4.west);
\draw[arr, minor!70!black] (p5.west) -- (open.east);
\node[chal, left=10mm of bus, anchor=east] {\emph{the blueprint's phase} \lean{bus}; seams
  \lean{Seam.commit} $\to$ \lean{Seam.bus}};
\node[chal, left=10mm of table, anchor=east] {\lean{table}; \lean{Seam.bus} $\to$
  \lean{Seam.table}};
\node[chal, left=10mm of pub, anchor=east] {\lean{pub}; \lean{Seam.table} $\to$
  \lean{Seam.pub}; \emph{built}};
\node[chal, left=10mm of flock, anchor=east] {\lean{flock}; \lean{Seam.pub} $\to$
  \lean{Seam.flock}};
\node[chal, left=10mm of open, anchor=east] {\lean{opening}; \lean{Seam.flock} $\to$
  \lean{Seam.done}};
\node[chal, left=10mm of commit, anchor=east] {\lean{commitDef}; \lean{M3Rel} $\to$
  \lean{Seam.commit}; \emph{built}};
\node[chal, left=10mm of setup, anchor=east] {not a phase: the compiled verifier's
  (Layer 12)};
\end{tikzpicture}
\caption{leanVM's proof system at the pin, as the deployed verifier runs it (centre), what each
phase adds to the claim pool (right), and the blueprint's phase and seams for it (left).}
\label{fig:phases}
\end{figure}

\section{What the blueprint makes of it}\label{sec:map-blueprint}

The blueprint models the protocol as a sequential composition of ArkLib \emph{oracle
reductions}, one per phase, in the ideal model where $q$ is an oracle and the challenges are
uniform (\cref{ch:libraries} introduces the vocabulary). It fixes, in the spine, what the
phases share: the abstract instance \lean{M3Instance}, the relation \lean{M3Holds}, the seam
relations, the interfaces \lean{Phase.Def}, \lean{Phase.Complete} and \lean{Phase.Security}, and
the composition with its two master theorems. It then specifies each phase (Layers 4 to 10),
the adaptor to leanISA (Layers 2 and 3), the compilation (Layers 11 and 12) and the target
theorem (Layer 13). \Cref{tab:map-where} says where each part is.

\begin{table}[h]
\centering\small
\begin{tabular}{@{}p{0.27\linewidth}p{0.30\linewidth}p{0.36\linewidth}@{}}
\toprule
Part & Blueprint & Code on \code{main} at \code{b435631} \\
\midrule
Fields, sampling, the oracle interface of $q$ & Layer 0 & \file{LeanerVM/Protocol/Field.lean}, \file{ToArkLib/Oracles.lean}; built \\
Hypercube tables, stacking, padding, the index and bytecode columns & Layer 1 & \file{ToCompPoly/*}, \file{Stack.lean}, \file{Padding.lean}, \file{ClaimWeights.lean}, \file{BlockClaims.lean}, \file{FixedColumns.lean}; built \\
The abstract instance, the relation, the seams, the interfaces, the composition, the master theorems & the section \enquote{The spine} & \file{Spine/\{Instance,Seams,Phase,Compose,Toy\}.lean}, \file{ToArkLib/\{Component,KnowledgeAppend,PassThrough,SendOracle,Refinement,GuardedVerdict,KeepOracles\}.lean}; built \\
Clean components as polynomials & Layer 2 & not built \\
The leanISA instance and the adaptor & Layer 3 & not built \\
The sumcheck and the batching component & Layer 4 & not built (an open pull request holds the honest round algebra) \\
Fingerprints, the grand product, GKR & Layer 5 & not built \\
The bus phase & Layer 6 & not built \\
The table sumcheck & Layer 7 & not built \\
The public-input phase & Layer 8 & \file{PublicInput.lean}; built, both halves proved \\
Flock and ring switching & Layer 9 & not built; owned by the Flock roadmap (issue \#3), which has not started \\
The claim pool and the opening & Layer 10 & not built \\
WHIR, Merkle trees, the parameters & Layer 11 & not built \\
Fiat--Shamir, the proof object, \lean{verify} & Layer 12 & not built \\
The target theorem & Layer 13 & not built \\
\bottomrule
\end{tabular}
\caption{Where each part of the proof system is specified and whether it is built.}
\label{tab:map-where}
\end{table}

Of the six phases, two are built: the commit phase, which is the spine's, and the public-input
phase. The four others, the generic components they rest on, the adaptor and the compilation
are specified only. The master theorems are therefore theorems about a \emph{bundle} of phases
supplied as a hypothesis; what they say, and what they leave to be trusted, is the subject of
\cref{ch:phases,ch:nonvacuity}.

\section{The reviewed revision and the one after it}\label{sec:map-revision}

The review's object is \code{main} at \code{b435631}. While the review was under way, the
owner upgraded \code{main} to Lean 4.34.1 (\code{144c5aa}, pull request \#61) and moved every
library pin: ArkLib to \code{7653a901}, CompPoly to \code{572f9973}, Clean to \code{42fe4b26},
VCVio to \code{a4232d08}. The leanVM pin did not move. The proof-system sources changed in
proofs and imports only (77 lines in, 76 out), so every statement this report discusses is the
same at both revisions; the two probability bridges of the spine were restated on VCVio's new
notation with the same relations and bounds. Two things did change in substance and are
recorded where they matter: ArkLib's new pin is the revision several dossiers examined as
\enquote{upstream}, and it lifts none of the admissions the blueprint's ledger names
(\cref{ch:libraries}); and CompPoly's $\K$ now gives natural-number literals their
characteristic-two meaning, $(2 : \K) = 0$, with encoded words written \lean{K.ofBits n}, so
that any probe or fixture using a literal other than $0$ or $1$ in $\K$ means something else at
the new pin (\cref{ch:probes}).
